%% =====================================================================
%% notation.tex　《高等代数讲义》本书记号约定
%% 用法：在 document.tex 的 \frontmatter 中、\tableofcontents 之前 %% =====================================================================
%% notation.tex　《高等代数讲义》本书记号约定
%% 用法：在 document.tex 的 \frontmatter 中、\tableofcontents 之前 %% =====================================================================
%% notation.tex　《高等代数讲义》本书记号约定
%% 用法：在 document.tex 的 \frontmatter 中、\tableofcontents 之前 %% =====================================================================
%% notation.tex　《高等代数讲义》本书记号约定
%% 用法：在 document.tex 的 \frontmatter 中、\tableofcontents 之前 \input{notation}.
%% 依赖：导言区定义的 \vv \zero \trans \spn \rk \imm \diag \tr \hB \hC.
%% =====================================================================

\chapter*{本书记号约定}
\addcontentsline{toc}{chapter}{本书记号约定}
\markboth{本书记号约定}{本书记号约定}

\paragraph*{关于书名.} 在我国高校的课程体系中，``高等代数''通常指数学类专业为期两个学期的基础课，除线性代数外还包括较系统的多项式理论和空间理论；面向非数学专业，如经济、管理、法学等专业开设的同类课程，则多称``线性代数''，以方程组、矩阵、向量空间、特征值与二次型为主. 本书不带标记的核心内容与后者的教学要求相当，按下文的 48 学时方案讲授即可. 书中另外收入了多项式、对偶空间与双线性型等传统上属于高等代数的内容，供 72、96 学时方案和有兴趣的读者选用，本书书名由此而来.

\section*{一、学时标记与教学方案}

本书内容明显多于一个学期所能讲授的分量. 为方便不同课时的课程选用，本书使用下列记号进行区分：

\begin{center}
	\begin{tabular}{cl}
		\toprule
		标记 & 含义 \\
		\midrule
		（无标记） & 核心内容，48 学时方案即讲授，72、96 学时方案也都讲授 \\
		\textdagger & 72 学时方案起讲授（96 学时方案也讲授） \\
		\textdaggerdbl & 仅在 96 学时方案中讲授 \\
		* & 选读内容，任何方案都不作要求，供有兴趣的读者自学 \\
		\bottomrule
	\end{tabular}
\end{center}

一节标题带有标记时，节内未另加标记的小节与它同属一个方案；小节自己的标记优先. 各章的建议学时如下表. 表中``—''表示该方案不讲这一章.

\begin{center}
	\begin{tabular}{llccc}
		\toprule
		章 & 内容 & 48 学时 & 72 学时 & 96 学时 \\
		\midrule
		第一章 & 线性方程组与消元法 & 3 & 3 & 3 \\
		第二章 & 向量空间 & 5 & 6 & 7 \\
		第三章 & 矩阵 & 6 & 7 & 8 \\
		第四章 & 线性空间与线性变换 & 5 & 7 & 9 \\
		第五章 & 行列式 & 4 & 5 & 6 \\
		第六章 & 线性方程组解的结构 & 4 & 5 & 7 \\
		第七章 & 多项式 & — & 4 & 6 \\
		第八章 & 特征值与特征向量 & 6 & 7 & 9 \\
		第九章 & $\lambda$-矩阵 & — & 4 & 6 \\
		第十章 & 内积空间 & 3 & 6 & 7 \\
		第十一章 & 矩阵分解 & 3 & 5 & 7 \\
		第十二章 & 酉空间 & — & — & 6 \\
		第十三章 & 二次型 & 5 & 7 & 8 \\
		第十四章 & *对偶空间与双线性型 & — & — & — \\
		\midrule
		& 复习与机动 & 4 & 6 & 7 \\
		& 合计 & 48 & 72 & 96 \\
		\bottomrule
	\end{tabular}
\end{center}

\paragraph*{使用说明.}
\begin{itemize}
	\item \textbf{48 学时方案}讲授各章不带标记的内容，依次覆盖：向量与矩阵运算、消元法；线性相关与秩；矩阵的运算与特殊矩阵；线性空间与线性映射；行列式；矩阵的秩；线性方程组的解；特征值与相似对角化；内积与欧氏空间；二次型. 
	\item 48 学时方案不讲第七章. 第八章用到的代数基本定理与 Vieta 定理，以不加证明的形式收在附录~\ref{app:complex}，所需的复数知识也在那里. 不带标记的内容若引用了带标记小节中的结论，48 学时方案只要求了解结论.
	\item \textbf{72 学时方案}讲授不带标记与带 \textdagger 的内容，大致相当于一学期、每周 4—5 学时. 第七章讲``一元多项式环''``带余除法''``根''``代数基本定理''与``最小多项式''；第八章不讲 Cayley–Hamilton 定理一节，因此对角化判据以定理~\ref{thm:diagcriterion} 为准；第九章 $\lambda$-矩阵整章讲授.
	\item \textbf{96 学时方案}讲完全部无星号的内容，包括第十二章酉空间，适合两个学期、每学期 48 学时的课程. 第十四章与各章星号小节可作为讨论班或期末报告的选题.
	\item 三个部分末尾的\textbf{综合项目}不占课内学时. 建议 48 学时方案选做第一部分的一个项目；72 学时方案在第一、二部分中各选一个；96 学时方案三个部分各选一个. 评分量规见``第一部分综合项目''.
	\item 章末习题的 A 组覆盖各方案的核心内容（涉及 \textdagger、\textdaggerdbl 小节的题目，未讲授时可以跳过）. A 组的答案、B 组的提示见附录.
\end{itemize}


\section*{二、栏目}

除了定义、定理与证明，各章还穿插着六类固定栏目，用不同颜色的方框标出. 它们各有分工，读者可以按需要取舍，但不宜一概跳过.

各章开头依次给出本章地位、建议学时、学习目标、前置知识与学习提示. 学习目标中注明``72 学时方案起''或``96 学时方案''的条目，对应正文中带 \textdagger 或 \textdaggerdbl 的内容；其余条目属于各方案共同的要求.

\textbf{思维动机}是正文最重要的补充. 它回答了另一个重要问题，即``为什么会想到这么做''：一个定义为什么这样写，一个证明的关键一步从何而来，一个结论在全书的位置在哪里. 初读时若只能挑一类栏目来读，编者推荐多留意这一类.

\textbf{专业连线}与\textbf{建模实验}即为数学建模的相关内容. 前者篇幅较短，指出一个概念在专业里的对应物，例如价目表是向量、投入产出表是矩阵；后者是完整的案例，依次走过``问题—假设—模型—求解—解读—局限''六步. 其中最后一步最容易被忽略，也最值得用心：每个模型都建立在简化之上，知道它在哪里失效，和会用它同样重要. 投入产出、物流网络、判例引证、法条检索等几个例子会跨章节反复出现，请留意这些问题是怎样随着工具的更新而被研究得更透彻的.

\textbf{历史回响}讲述概念的来历：谁在什么问题的推动下提出了它，又经过了怎样的曲折. 栏目中的史实均注明出处，见书末参考文献. \textbf{常见误区}收集初学者最容易犯的错误，多数以反例的形式给出，适合在做习题之前进行阅读.

\textbf{计算实践}是使用 Python 进行实践的栏目. 栏目中会布置数个任务供读者进行探究，并附检查点用于检验结果. 除此以外，本栏目还有介绍实际的数值计算和计算机求解中的算法的作用.

\section*{三、编号与引用}
\begin{itemize}
	\item 定义、定理、引理、推论、命题、例题按节编号，例如``定理 2.4.3''是第二章第四节的第 3 个定理.
	\item 章末习题分为 A、B、C三组，每组各自编号. 正文中``第三章章末习题 B16''指第三章章末习题 B 组第 16 题. 正文各节中零星的``练习''按出现顺序编号，与章末习题不同.
	\item 证明结束用 $\blacksquare$ 标出.
\end{itemize}

\section*{四、本书数学符号一览}

\textbf{数与集合}

\begin{center}
	\begin{longtable}{l>{\raggedright\arraybackslash}p{0.6\linewidth}}
		\toprule
		记号 & 含义 \\
		\midrule
		\endhead
		$\mathbb{N},\ \mathbb{Z},\ \mathbb{Q},\ \mathbb{R},\ \mathbb{C}$ & 自然数集、整数集、有理数集、实数集、复数集 \\
		$\mathbb{F}$ & 一般的数域 \\
		$A \subset B$ & $A$ 是 $B$ 的子集，等价于 $B \supset A$ \\
		$\mathrm{i}$ & 虚数单位，$\mathrm i^2=-1$ \\
		$\bar z$ & 复数 $z$ 的共轭 \\
		$|z|$ & 复数 $z$ 的模 \\
		$\mathrm e^{\mathrm i\theta}$ & $\cos\theta+\mathrm i\sin\theta$（见附录``复数''一节） \\
		$\omega$ & 三次单位根 $\mathrm e^{2\pi\mathrm i/3}$ \\
		$\mathbb{F}_2$ & （取决于上下文）只含 $0,1$、按模 $2$ 运算的数系 \\
		$\sum_{i=1}^n a_i$，$\prod_{i=1}^na_i$ & 求和、求积（见附录``预备知识''） \\
		$\binom nk$ & 组合数 \\
		$\deg f$ & 多项式 $f$ 的次数 \\
		$g\mid f$ & 多项式（或整数）$g$ 整除 $f$ \\
		$\gcd(f,g)$ & 首一的最大公因式 \\
		$f:A\to B$，$x\mapsto f(x)$ & 映射及其对应法则 \\
		$g\circ f$ & 映射的复合，先 $f$ 后 $g$ \\
		$\mathbb{F}[x]$，$\mathbb{F}[x]_n$ & $\mathbb{F}$ 上的一元多项式全体；次数小于 $n$ 的多项式（含零多项式）全体 \\
		$C[a,b]$ & 闭区间 $[a,b]$ 上的实值连续函数全体 \\
		\bottomrule
	\end{longtable}
\end{center}
\newpage
\textbf{向量与矩阵}

\begin{center}
	\begin{longtable}{l>{\raggedright\arraybackslash}p{0.6\linewidth}}
		\toprule
		记号 & 含义 \\
		\midrule
		\endhead
		$\vv\alpha,\ \vv\beta,\ \vv x,\ \vv A$ & 向量与矩阵一律用加粗斜体 \\
		$\zero$ & 零向量 \\
		$\vv O$ & 零矩阵 \\
		$\vv I$，$\vv I_n$ & 单位矩阵、$n$ 阶单位矩阵 \\
		$\vv e_i$ & 标准单位向量：第 $i$ 个分量为 $1$，其余为 $0$ \\
		$\vv 1$ & 分量全为 $1$ 的列向量 \\
		$\mathbb{F}^n$ & $\mathbb{F}$ 上的 $n$ 维列向量全体；本书``向量''默认指列向量 \\
		$M_{m\times n}(\mathbb{F})$，$M_n(\mathbb{F})$ & $m\times n$ 矩阵全体；$n$ 阶方阵全体 \\
		$a_{ij}$，$(a_{ij})_{m\times n}$ & 矩阵的 $(i,j)$ 元；以 $a_{ij}$ 为元素的矩阵 \\
		$\vv A^\trans$ & 转置 \\
		$\vv A^{-1}$ & 逆矩阵 \\
		$\vv A^*$ & 伴随矩阵（第五章）；第十四章中 $\mathscr T^*$ 另指对偶映射 \\
		$(\vv A\mid\vv b)$ & 增广矩阵 \\
		$\diag(\lambda_1,\ldots,\lambda_n)$ & 对角矩阵 \\
		$r_i\leftrightarrow r_j$，$kr_i$，$r_i+cr_j$ & 三类初等行变换（列变换用 $c$ 代替 $r$） \\
		$\vv E(i,j)$，$\vv E(i(k))$，$\vv E(i,j(c))$ & 三类初等矩阵 \\
		$\rk\vv A$ & 秩 \\
		$\det\vv A$，$|\vv A|$ & 行列式 \\
		$\tr\vv A$ & 迹，即主对角元之和 \\
		$M_{ij}$，$A_{ij}$ & 余子式、代数余子式 \\
		$\tau(i_1i_2\cdots i_n)$ & 排列的逆序数 \\
		$f_{\vv A}(\lambda)$，$m_{\vv A}(\lambda)$ & 特征多项式 $\det(\lambda\vv I-\vv A)$，最小多项式 \\
		$\vv J_k(\lambda)$ & $k$ 阶 Jordan 块 \\
		$\sigma_1\ge\sigma_2\ge\cdots$ & 奇异值 \\
		$\kappa(\vv A)$ & 条件数 $\sigma_1/\sigma_n$ \\
		$\|\vv A\|_F$ & Frobenius 范数 $\sqrt{\sum a_{ij}^2}$ \\
				$\|\vv A\|_2$ & 谱范数，即最大奇异值 $\sigma_1$ \\
		$\vv A^+$ & 伪逆 \\
		$\mathrm e^{t\vv A}$ & 矩阵指数 \\
		$\vv R(\theta)$，$\vv F(\theta)$ & 平面上转角 $\theta$ 的旋转矩阵；关于倾角 $\theta/2$ 的直线的反射矩阵 \\
		$\vv H$ & Householder 矩阵 $\vv I-2\vv u\vv u^\trans$ \\
		$\vv A\otimes\vv B$，$\vv A\circ\vv B$ & Kronecker 积；逐元素乘积（仅见于 C 组习题） \\
		\bottomrule
	\end{longtable}
\end{center}
\newpage
\textbf{空间与映射}

\begin{center}
	\begin{longtable}{l>{\raggedright\arraybackslash}p{0.6\linewidth}}
		\toprule
		记号 & 含义 \\
		\midrule
		\endhead
		$\spn(\vv\alpha_1,\ldots,\vv\alpha_s)$ & 向量组张成的子空间 \\
		$\dim V$，$\dim_{\mathbb{F}}V$ & 维数（需要强调数域时写出 $\mathbb{F}$） \\
		$W_1+W_2$，$W_1\oplus W_2$ & 子空间的和、直和 \\
		$V\times W$ & 向量空间的直积 \\
		$V\cong U$ & 同构 \\
		$\mathscr T,\ \mathscr S,\ \mathscr D$ & 线性映射与线性变换用花体字母；$\mathscr I$ 是恒等变换，$\mathscr O$ 是零变换 \\
		$\ker\mathscr T$，$\imm\mathscr T$ & 核、像；对矩阵 $\vv A$，$\ker\vv A$、$\imm\vv A$ 指 $\vv x\mapsto\vv A\vv x$ 的核与像 \\
		$\mathscr T|_W$ & $\mathscr T$ 在不变子空间 $W$ 上的限制 \\
		$L(V,W)$ & $V$ 到 $W$ 的全体线性映射 \\
		$V_\lambda$ & 属于特征值 $\lambda$ 的特征子空间 \\
		$\langle\vv\alpha,\vv\beta\rangle$ & 内积；$\mathbb{R}^n$ 中默认为标准内积 $\vv\alpha^\trans\vv\beta$ \\
		$\|\vv\alpha\|$ & 范数（长度） \\
		$\vv\alpha\perp\vv\beta$ & 正交 \\
		$W^\perp$ & 正交补 \\
		$\vv\alpha_W$ & $\vv\alpha$ 在子空间 $W$ 上的正交投影 \\
		$\vv x_0+W$ & 线性流形（仿射子空间） \\
		$V^*$，$W^0$，$\mathscr T^*$ & 对偶空间、零化子、对偶映射（第十四章） \\
		$\hat{\vv\alpha}$ & $\vv\alpha$ 在 $V^{**}$ 中的像：$f\mapsto f(\vv\alpha)$（第十四章） \\
		$\vv x\ge\zero$ & 向量的每个分量都非负 \\
		\bottomrule
	\end{longtable}
\end{center}

\begin{note}
	同一个字母在不同章节中可能另有所指，例如 $\vv P$ 既表示过渡矩阵，也表示转移矩阵与投影矩阵；$\vv\Sigma$ 在第十一章是奇异值构成的矩阵，在第十三章是协方差矩阵. 这时以该章的说明为准.
\end{note}
.
%% 依赖：导言区定义的 \vv \zero \trans \spn \rk \imm \diag \tr \hB \hC.
%% =====================================================================

\chapter*{本书记号约定}
\addcontentsline{toc}{chapter}{本书记号约定}
\markboth{本书记号约定}{本书记号约定}

\paragraph*{关于书名.} 在我国高校的课程体系中，``高等代数''通常指数学类专业为期两个学期的基础课，除线性代数外还包括较系统的多项式理论和空间理论；面向非数学专业，如经济、管理、法学等专业开设的同类课程，则多称``线性代数''，以方程组、矩阵、向量空间、特征值与二次型为主. 本书不带标记的核心内容与后者的教学要求相当，按下文的 48 学时方案讲授即可. 书中另外收入了多项式、对偶空间与双线性型等传统上属于高等代数的内容，供 72、96 学时方案和有兴趣的读者选用，本书书名由此而来.

\section*{一、学时标记与教学方案}

本书内容明显多于一个学期所能讲授的分量. 为方便不同课时的课程选用，本书使用下列记号进行区分：

\begin{center}
	\begin{tabular}{cl}
		\toprule
		标记 & 含义 \\
		\midrule
		（无标记） & 核心内容，48 学时方案即讲授，72、96 学时方案也都讲授 \\
		\textdagger & 72 学时方案起讲授（96 学时方案也讲授） \\
		\textdaggerdbl & 仅在 96 学时方案中讲授 \\
		* & 选读内容，任何方案都不作要求，供有兴趣的读者自学 \\
		\bottomrule
	\end{tabular}
\end{center}

一节标题带有标记时，节内未另加标记的小节与它同属一个方案；小节自己的标记优先. 各章的建议学时如下表. 表中``—''表示该方案不讲这一章.

\begin{center}
	\begin{tabular}{llccc}
		\toprule
		章 & 内容 & 48 学时 & 72 学时 & 96 学时 \\
		\midrule
		第一章 & 线性方程组与消元法 & 3 & 3 & 3 \\
		第二章 & 向量空间 & 5 & 6 & 7 \\
		第三章 & 矩阵 & 6 & 7 & 8 \\
		第四章 & 线性空间与线性变换 & 5 & 7 & 9 \\
		第五章 & 行列式 & 4 & 5 & 6 \\
		第六章 & 线性方程组解的结构 & 4 & 5 & 7 \\
		第七章 & 多项式 & — & 4 & 6 \\
		第八章 & 特征值与特征向量 & 6 & 7 & 9 \\
		第九章 & $\lambda$-矩阵 & — & 4 & 6 \\
		第十章 & 内积空间 & 3 & 6 & 7 \\
		第十一章 & 矩阵分解 & 3 & 5 & 7 \\
		第十二章 & 酉空间 & — & — & 6 \\
		第十三章 & 二次型 & 5 & 7 & 8 \\
		第十四章 & *对偶空间与双线性型 & — & — & — \\
		\midrule
		& 复习与机动 & 4 & 6 & 7 \\
		& 合计 & 48 & 72 & 96 \\
		\bottomrule
	\end{tabular}
\end{center}

\paragraph*{使用说明.}
\begin{itemize}
	\item \textbf{48 学时方案}讲授各章不带标记的内容，依次覆盖：向量与矩阵运算、消元法；线性相关与秩；矩阵的运算与特殊矩阵；线性空间与线性映射；行列式；矩阵的秩；线性方程组的解；特征值与相似对角化；内积与欧氏空间；二次型. 
	\item 48 学时方案不讲第七章. 第八章用到的代数基本定理与 Vieta 定理，以不加证明的形式收在附录~\ref{app:complex}，所需的复数知识也在那里. 不带标记的内容若引用了带标记小节中的结论，48 学时方案只要求了解结论.
	\item \textbf{72 学时方案}讲授不带标记与带 \textdagger 的内容，大致相当于一学期、每周 4—5 学时. 第七章讲``一元多项式环''``带余除法''``根''``代数基本定理''与``最小多项式''；第八章不讲 Cayley–Hamilton 定理一节，因此对角化判据以定理~\ref{thm:diagcriterion} 为准；第九章 $\lambda$-矩阵整章讲授.
	\item \textbf{96 学时方案}讲完全部无星号的内容，包括第十二章酉空间，适合两个学期、每学期 48 学时的课程. 第十四章与各章星号小节可作为讨论班或期末报告的选题.
	\item 三个部分末尾的\textbf{综合项目}不占课内学时. 建议 48 学时方案选做第一部分的一个项目；72 学时方案在第一、二部分中各选一个；96 学时方案三个部分各选一个. 评分量规见``第一部分综合项目''.
	\item 章末习题的 A 组覆盖各方案的核心内容（涉及 \textdagger、\textdaggerdbl 小节的题目，未讲授时可以跳过）. A 组的答案、B 组的提示见附录.
\end{itemize}


\section*{二、栏目}

除了定义、定理与证明，各章还穿插着六类固定栏目，用不同颜色的方框标出. 它们各有分工，读者可以按需要取舍，但不宜一概跳过.

各章开头依次给出本章地位、建议学时、学习目标、前置知识与学习提示. 学习目标中注明``72 学时方案起''或``96 学时方案''的条目，对应正文中带 \textdagger 或 \textdaggerdbl 的内容；其余条目属于各方案共同的要求.

\textbf{思维动机}是正文最重要的补充. 它回答了另一个重要问题，即``为什么会想到这么做''：一个定义为什么这样写，一个证明的关键一步从何而来，一个结论在全书的位置在哪里. 初读时若只能挑一类栏目来读，编者推荐多留意这一类.

\textbf{专业连线}与\textbf{建模实验}即为数学建模的相关内容. 前者篇幅较短，指出一个概念在专业里的对应物，例如价目表是向量、投入产出表是矩阵；后者是完整的案例，依次走过``问题—假设—模型—求解—解读—局限''六步. 其中最后一步最容易被忽略，也最值得用心：每个模型都建立在简化之上，知道它在哪里失效，和会用它同样重要. 投入产出、物流网络、判例引证、法条检索等几个例子会跨章节反复出现，请留意这些问题是怎样随着工具的更新而被研究得更透彻的.

\textbf{历史回响}讲述概念的来历：谁在什么问题的推动下提出了它，又经过了怎样的曲折. 栏目中的史实均注明出处，见书末参考文献. \textbf{常见误区}收集初学者最容易犯的错误，多数以反例的形式给出，适合在做习题之前进行阅读.

\textbf{计算实践}是使用 Python 进行实践的栏目. 栏目中会布置数个任务供读者进行探究，并附检查点用于检验结果. 除此以外，本栏目还有介绍实际的数值计算和计算机求解中的算法的作用.

\section*{三、编号与引用}
\begin{itemize}
	\item 定义、定理、引理、推论、命题、例题按节编号，例如``定理 2.4.3''是第二章第四节的第 3 个定理.
	\item 章末习题分为 A、B、C三组，每组各自编号. 正文中``第三章章末习题 B16''指第三章章末习题 B 组第 16 题. 正文各节中零星的``练习''按出现顺序编号，与章末习题不同.
	\item 证明结束用 $\blacksquare$ 标出.
\end{itemize}

\section*{四、本书数学符号一览}

\textbf{数与集合}

\begin{center}
	\begin{longtable}{l>{\raggedright\arraybackslash}p{0.6\linewidth}}
		\toprule
		记号 & 含义 \\
		\midrule
		\endhead
		$\mathbb{N},\ \mathbb{Z},\ \mathbb{Q},\ \mathbb{R},\ \mathbb{C}$ & 自然数集、整数集、有理数集、实数集、复数集 \\
		$\mathbb{F}$ & 一般的数域 \\
		$A \subset B$ & $A$ 是 $B$ 的子集，等价于 $B \supset A$ \\
		$\mathrm{i}$ & 虚数单位，$\mathrm i^2=-1$ \\
		$\bar z$ & 复数 $z$ 的共轭 \\
		$|z|$ & 复数 $z$ 的模 \\
		$\mathrm e^{\mathrm i\theta}$ & $\cos\theta+\mathrm i\sin\theta$（见附录``复数''一节） \\
		$\omega$ & 三次单位根 $\mathrm e^{2\pi\mathrm i/3}$ \\
		$\mathbb{F}_2$ & （取决于上下文）只含 $0,1$、按模 $2$ 运算的数系 \\
		$\sum_{i=1}^n a_i$，$\prod_{i=1}^na_i$ & 求和、求积（见附录``预备知识''） \\
		$\binom nk$ & 组合数 \\
		$\deg f$ & 多项式 $f$ 的次数 \\
		$g\mid f$ & 多项式（或整数）$g$ 整除 $f$ \\
		$\gcd(f,g)$ & 首一的最大公因式 \\
		$f:A\to B$，$x\mapsto f(x)$ & 映射及其对应法则 \\
		$g\circ f$ & 映射的复合，先 $f$ 后 $g$ \\
		$\mathbb{F}[x]$，$\mathbb{F}[x]_n$ & $\mathbb{F}$ 上的一元多项式全体；次数小于 $n$ 的多项式（含零多项式）全体 \\
		$C[a,b]$ & 闭区间 $[a,b]$ 上的实值连续函数全体 \\
		\bottomrule
	\end{longtable}
\end{center}
\newpage
\textbf{向量与矩阵}

\begin{center}
	\begin{longtable}{l>{\raggedright\arraybackslash}p{0.6\linewidth}}
		\toprule
		记号 & 含义 \\
		\midrule
		\endhead
		$\vv\alpha,\ \vv\beta,\ \vv x,\ \vv A$ & 向量与矩阵一律用加粗斜体 \\
		$\zero$ & 零向量 \\
		$\vv O$ & 零矩阵 \\
		$\vv I$，$\vv I_n$ & 单位矩阵、$n$ 阶单位矩阵 \\
		$\vv e_i$ & 标准单位向量：第 $i$ 个分量为 $1$，其余为 $0$ \\
		$\vv 1$ & 分量全为 $1$ 的列向量 \\
		$\mathbb{F}^n$ & $\mathbb{F}$ 上的 $n$ 维列向量全体；本书``向量''默认指列向量 \\
		$M_{m\times n}(\mathbb{F})$，$M_n(\mathbb{F})$ & $m\times n$ 矩阵全体；$n$ 阶方阵全体 \\
		$a_{ij}$，$(a_{ij})_{m\times n}$ & 矩阵的 $(i,j)$ 元；以 $a_{ij}$ 为元素的矩阵 \\
		$\vv A^\trans$ & 转置 \\
		$\vv A^{-1}$ & 逆矩阵 \\
		$\vv A^*$ & 伴随矩阵（第五章）；第十四章中 $\mathscr T^*$ 另指对偶映射 \\
		$(\vv A\mid\vv b)$ & 增广矩阵 \\
		$\diag(\lambda_1,\ldots,\lambda_n)$ & 对角矩阵 \\
		$r_i\leftrightarrow r_j$，$kr_i$，$r_i+cr_j$ & 三类初等行变换（列变换用 $c$ 代替 $r$） \\
		$\vv E(i,j)$，$\vv E(i(k))$，$\vv E(i,j(c))$ & 三类初等矩阵 \\
		$\rk\vv A$ & 秩 \\
		$\det\vv A$，$|\vv A|$ & 行列式 \\
		$\tr\vv A$ & 迹，即主对角元之和 \\
		$M_{ij}$，$A_{ij}$ & 余子式、代数余子式 \\
		$\tau(i_1i_2\cdots i_n)$ & 排列的逆序数 \\
		$f_{\vv A}(\lambda)$，$m_{\vv A}(\lambda)$ & 特征多项式 $\det(\lambda\vv I-\vv A)$，最小多项式 \\
		$\vv J_k(\lambda)$ & $k$ 阶 Jordan 块 \\
		$\sigma_1\ge\sigma_2\ge\cdots$ & 奇异值 \\
		$\kappa(\vv A)$ & 条件数 $\sigma_1/\sigma_n$ \\
		$\|\vv A\|_F$ & Frobenius 范数 $\sqrt{\sum a_{ij}^2}$ \\
				$\|\vv A\|_2$ & 谱范数，即最大奇异值 $\sigma_1$ \\
		$\vv A^+$ & 伪逆 \\
		$\mathrm e^{t\vv A}$ & 矩阵指数 \\
		$\vv R(\theta)$，$\vv F(\theta)$ & 平面上转角 $\theta$ 的旋转矩阵；关于倾角 $\theta/2$ 的直线的反射矩阵 \\
		$\vv H$ & Householder 矩阵 $\vv I-2\vv u\vv u^\trans$ \\
		$\vv A\otimes\vv B$，$\vv A\circ\vv B$ & Kronecker 积；逐元素乘积（仅见于 C 组习题） \\
		\bottomrule
	\end{longtable}
\end{center}
\newpage
\textbf{空间与映射}

\begin{center}
	\begin{longtable}{l>{\raggedright\arraybackslash}p{0.6\linewidth}}
		\toprule
		记号 & 含义 \\
		\midrule
		\endhead
		$\spn(\vv\alpha_1,\ldots,\vv\alpha_s)$ & 向量组张成的子空间 \\
		$\dim V$，$\dim_{\mathbb{F}}V$ & 维数（需要强调数域时写出 $\mathbb{F}$） \\
		$W_1+W_2$，$W_1\oplus W_2$ & 子空间的和、直和 \\
		$V\times W$ & 向量空间的直积 \\
		$V\cong U$ & 同构 \\
		$\mathscr T,\ \mathscr S,\ \mathscr D$ & 线性映射与线性变换用花体字母；$\mathscr I$ 是恒等变换，$\mathscr O$ 是零变换 \\
		$\ker\mathscr T$，$\imm\mathscr T$ & 核、像；对矩阵 $\vv A$，$\ker\vv A$、$\imm\vv A$ 指 $\vv x\mapsto\vv A\vv x$ 的核与像 \\
		$\mathscr T|_W$ & $\mathscr T$ 在不变子空间 $W$ 上的限制 \\
		$L(V,W)$ & $V$ 到 $W$ 的全体线性映射 \\
		$V_\lambda$ & 属于特征值 $\lambda$ 的特征子空间 \\
		$\langle\vv\alpha,\vv\beta\rangle$ & 内积；$\mathbb{R}^n$ 中默认为标准内积 $\vv\alpha^\trans\vv\beta$ \\
		$\|\vv\alpha\|$ & 范数（长度） \\
		$\vv\alpha\perp\vv\beta$ & 正交 \\
		$W^\perp$ & 正交补 \\
		$\vv\alpha_W$ & $\vv\alpha$ 在子空间 $W$ 上的正交投影 \\
		$\vv x_0+W$ & 线性流形（仿射子空间） \\
		$V^*$，$W^0$，$\mathscr T^*$ & 对偶空间、零化子、对偶映射（第十四章） \\
		$\hat{\vv\alpha}$ & $\vv\alpha$ 在 $V^{**}$ 中的像：$f\mapsto f(\vv\alpha)$（第十四章） \\
		$\vv x\ge\zero$ & 向量的每个分量都非负 \\
		\bottomrule
	\end{longtable}
\end{center}

\begin{note}
	同一个字母在不同章节中可能另有所指，例如 $\vv P$ 既表示过渡矩阵，也表示转移矩阵与投影矩阵；$\vv\Sigma$ 在第十一章是奇异值构成的矩阵，在第十三章是协方差矩阵. 这时以该章的说明为准.
\end{note}
.
%% 依赖：导言区定义的 \vv \zero \trans \spn \rk \imm \diag \tr \hB \hC.
%% =====================================================================

\chapter*{本书记号约定}
\addcontentsline{toc}{chapter}{本书记号约定}
\markboth{本书记号约定}{本书记号约定}

\paragraph*{关于书名.} 在我国高校的课程体系中，``高等代数''通常指数学类专业为期两个学期的基础课，除线性代数外还包括较系统的多项式理论和空间理论；面向非数学专业，如经济、管理、法学等专业开设的同类课程，则多称``线性代数''，以方程组、矩阵、向量空间、特征值与二次型为主. 本书不带标记的核心内容与后者的教学要求相当，按下文的 48 学时方案讲授即可. 书中另外收入了多项式、对偶空间与双线性型等传统上属于高等代数的内容，供 72、96 学时方案和有兴趣的读者选用，本书书名由此而来.

\section*{一、学时标记与教学方案}

本书内容明显多于一个学期所能讲授的分量. 为方便不同课时的课程选用，本书使用下列记号进行区分：

\begin{center}
	\begin{tabular}{cl}
		\toprule
		标记 & 含义 \\
		\midrule
		（无标记） & 核心内容，48 学时方案即讲授，72、96 学时方案也都讲授 \\
		\textdagger & 72 学时方案起讲授（96 学时方案也讲授） \\
		\textdaggerdbl & 仅在 96 学时方案中讲授 \\
		* & 选读内容，任何方案都不作要求，供有兴趣的读者自学 \\
		\bottomrule
	\end{tabular}
\end{center}

一节标题带有标记时，节内未另加标记的小节与它同属一个方案；小节自己的标记优先. 各章的建议学时如下表. 表中``—''表示该方案不讲这一章.

\begin{center}
	\begin{tabular}{llccc}
		\toprule
		章 & 内容 & 48 学时 & 72 学时 & 96 学时 \\
		\midrule
		第一章 & 线性方程组与消元法 & 3 & 3 & 3 \\
		第二章 & 向量空间 & 5 & 6 & 7 \\
		第三章 & 矩阵 & 6 & 7 & 8 \\
		第四章 & 线性空间与线性变换 & 5 & 7 & 9 \\
		第五章 & 行列式 & 4 & 5 & 6 \\
		第六章 & 线性方程组解的结构 & 4 & 5 & 7 \\
		第七章 & 多项式 & — & 4 & 6 \\
		第八章 & 特征值与特征向量 & 6 & 7 & 9 \\
		第九章 & $\lambda$-矩阵 & — & 4 & 6 \\
		第十章 & 内积空间 & 3 & 6 & 7 \\
		第十一章 & 矩阵分解 & 3 & 5 & 7 \\
		第十二章 & 酉空间 & — & — & 6 \\
		第十三章 & 二次型 & 5 & 7 & 8 \\
		第十四章 & *对偶空间与双线性型 & — & — & — \\
		\midrule
		& 复习与机动 & 4 & 6 & 7 \\
		& 合计 & 48 & 72 & 96 \\
		\bottomrule
	\end{tabular}
\end{center}

\paragraph*{使用说明.}
\begin{itemize}
	\item \textbf{48 学时方案}讲授各章不带标记的内容，依次覆盖：向量与矩阵运算、消元法；线性相关与秩；矩阵的运算与特殊矩阵；线性空间与线性映射；行列式；矩阵的秩；线性方程组的解；特征值与相似对角化；内积与欧氏空间；二次型. 
	\item 48 学时方案不讲第七章. 第八章用到的代数基本定理与 Vieta 定理，以不加证明的形式收在附录~\ref{app:complex}，所需的复数知识也在那里. 不带标记的内容若引用了带标记小节中的结论，48 学时方案只要求了解结论.
	\item \textbf{72 学时方案}讲授不带标记与带 \textdagger 的内容，大致相当于一学期、每周 4—5 学时. 第七章讲``一元多项式环''``带余除法''``根''``代数基本定理''与``最小多项式''；第八章不讲 Cayley–Hamilton 定理一节，因此对角化判据以定理~\ref{thm:diagcriterion} 为准；第九章 $\lambda$-矩阵整章讲授.
	\item \textbf{96 学时方案}讲完全部无星号的内容，包括第十二章酉空间，适合两个学期、每学期 48 学时的课程. 第十四章与各章星号小节可作为讨论班或期末报告的选题.
	\item 三个部分末尾的\textbf{综合项目}不占课内学时. 建议 48 学时方案选做第一部分的一个项目；72 学时方案在第一、二部分中各选一个；96 学时方案三个部分各选一个. 评分量规见``第一部分综合项目''.
	\item 章末习题的 A 组覆盖各方案的核心内容（涉及 \textdagger、\textdaggerdbl 小节的题目，未讲授时可以跳过）. A 组的答案、B 组的提示见附录.
\end{itemize}


\section*{二、栏目}

除了定义、定理与证明，各章还穿插着六类固定栏目，用不同颜色的方框标出. 它们各有分工，读者可以按需要取舍，但不宜一概跳过.

各章开头依次给出本章地位、建议学时、学习目标、前置知识与学习提示. 学习目标中注明``72 学时方案起''或``96 学时方案''的条目，对应正文中带 \textdagger 或 \textdaggerdbl 的内容；其余条目属于各方案共同的要求.

\textbf{思维动机}是正文最重要的补充. 它回答了另一个重要问题，即``为什么会想到这么做''：一个定义为什么这样写，一个证明的关键一步从何而来，一个结论在全书的位置在哪里. 初读时若只能挑一类栏目来读，编者推荐多留意这一类.

\textbf{专业连线}与\textbf{建模实验}即为数学建模的相关内容. 前者篇幅较短，指出一个概念在专业里的对应物，例如价目表是向量、投入产出表是矩阵；后者是完整的案例，依次走过``问题—假设—模型—求解—解读—局限''六步. 其中最后一步最容易被忽略，也最值得用心：每个模型都建立在简化之上，知道它在哪里失效，和会用它同样重要. 投入产出、物流网络、判例引证、法条检索等几个例子会跨章节反复出现，请留意这些问题是怎样随着工具的更新而被研究得更透彻的.

\textbf{历史回响}讲述概念的来历：谁在什么问题的推动下提出了它，又经过了怎样的曲折. 栏目中的史实均注明出处，见书末参考文献. \textbf{常见误区}收集初学者最容易犯的错误，多数以反例的形式给出，适合在做习题之前进行阅读.

\textbf{计算实践}是使用 Python 进行实践的栏目. 栏目中会布置数个任务供读者进行探究，并附检查点用于检验结果. 除此以外，本栏目还有介绍实际的数值计算和计算机求解中的算法的作用.

\section*{三、编号与引用}
\begin{itemize}
	\item 定义、定理、引理、推论、命题、例题按节编号，例如``定理 2.4.3''是第二章第四节的第 3 个定理.
	\item 章末习题分为 A、B、C三组，每组各自编号. 正文中``第三章章末习题 B16''指第三章章末习题 B 组第 16 题. 正文各节中零星的``练习''按出现顺序编号，与章末习题不同.
	\item 证明结束用 $\blacksquare$ 标出.
\end{itemize}

\section*{四、本书数学符号一览}

\textbf{数与集合}

\begin{center}
	\begin{longtable}{l>{\raggedright\arraybackslash}p{0.6\linewidth}}
		\toprule
		记号 & 含义 \\
		\midrule
		\endhead
		$\mathbb{N},\ \mathbb{Z},\ \mathbb{Q},\ \mathbb{R},\ \mathbb{C}$ & 自然数集、整数集、有理数集、实数集、复数集 \\
		$\mathbb{F}$ & 一般的数域 \\
		$A \subset B$ & $A$ 是 $B$ 的子集，等价于 $B \supset A$ \\
		$\mathrm{i}$ & 虚数单位，$\mathrm i^2=-1$ \\
		$\bar z$ & 复数 $z$ 的共轭 \\
		$|z|$ & 复数 $z$ 的模 \\
		$\mathrm e^{\mathrm i\theta}$ & $\cos\theta+\mathrm i\sin\theta$（见附录``复数''一节） \\
		$\omega$ & 三次单位根 $\mathrm e^{2\pi\mathrm i/3}$ \\
		$\mathbb{F}_2$ & （取决于上下文）只含 $0,1$、按模 $2$ 运算的数系 \\
		$\sum_{i=1}^n a_i$，$\prod_{i=1}^na_i$ & 求和、求积（见附录``预备知识''） \\
		$\binom nk$ & 组合数 \\
		$\deg f$ & 多项式 $f$ 的次数 \\
		$g\mid f$ & 多项式（或整数）$g$ 整除 $f$ \\
		$\gcd(f,g)$ & 首一的最大公因式 \\
		$f:A\to B$，$x\mapsto f(x)$ & 映射及其对应法则 \\
		$g\circ f$ & 映射的复合，先 $f$ 后 $g$ \\
		$\mathbb{F}[x]$，$\mathbb{F}[x]_n$ & $\mathbb{F}$ 上的一元多项式全体；次数小于 $n$ 的多项式（含零多项式）全体 \\
		$C[a,b]$ & 闭区间 $[a,b]$ 上的实值连续函数全体 \\
		\bottomrule
	\end{longtable}
\end{center}
\newpage
\textbf{向量与矩阵}

\begin{center}
	\begin{longtable}{l>{\raggedright\arraybackslash}p{0.6\linewidth}}
		\toprule
		记号 & 含义 \\
		\midrule
		\endhead
		$\vv\alpha,\ \vv\beta,\ \vv x,\ \vv A$ & 向量与矩阵一律用加粗斜体 \\
		$\zero$ & 零向量 \\
		$\vv O$ & 零矩阵 \\
		$\vv I$，$\vv I_n$ & 单位矩阵、$n$ 阶单位矩阵 \\
		$\vv e_i$ & 标准单位向量：第 $i$ 个分量为 $1$，其余为 $0$ \\
		$\vv 1$ & 分量全为 $1$ 的列向量 \\
		$\mathbb{F}^n$ & $\mathbb{F}$ 上的 $n$ 维列向量全体；本书``向量''默认指列向量 \\
		$M_{m\times n}(\mathbb{F})$，$M_n(\mathbb{F})$ & $m\times n$ 矩阵全体；$n$ 阶方阵全体 \\
		$a_{ij}$，$(a_{ij})_{m\times n}$ & 矩阵的 $(i,j)$ 元；以 $a_{ij}$ 为元素的矩阵 \\
		$\vv A^\trans$ & 转置 \\
		$\vv A^{-1}$ & 逆矩阵 \\
		$\vv A^*$ & 伴随矩阵（第五章）；第十四章中 $\mathscr T^*$ 另指对偶映射 \\
		$(\vv A\mid\vv b)$ & 增广矩阵 \\
		$\diag(\lambda_1,\ldots,\lambda_n)$ & 对角矩阵 \\
		$r_i\leftrightarrow r_j$，$kr_i$，$r_i+cr_j$ & 三类初等行变换（列变换用 $c$ 代替 $r$） \\
		$\vv E(i,j)$，$\vv E(i(k))$，$\vv E(i,j(c))$ & 三类初等矩阵 \\
		$\rk\vv A$ & 秩 \\
		$\det\vv A$，$|\vv A|$ & 行列式 \\
		$\tr\vv A$ & 迹，即主对角元之和 \\
		$M_{ij}$，$A_{ij}$ & 余子式、代数余子式 \\
		$\tau(i_1i_2\cdots i_n)$ & 排列的逆序数 \\
		$f_{\vv A}(\lambda)$，$m_{\vv A}(\lambda)$ & 特征多项式 $\det(\lambda\vv I-\vv A)$，最小多项式 \\
		$\vv J_k(\lambda)$ & $k$ 阶 Jordan 块 \\
		$\sigma_1\ge\sigma_2\ge\cdots$ & 奇异值 \\
		$\kappa(\vv A)$ & 条件数 $\sigma_1/\sigma_n$ \\
		$\|\vv A\|_F$ & Frobenius 范数 $\sqrt{\sum a_{ij}^2}$ \\
				$\|\vv A\|_2$ & 谱范数，即最大奇异值 $\sigma_1$ \\
		$\vv A^+$ & 伪逆 \\
		$\mathrm e^{t\vv A}$ & 矩阵指数 \\
		$\vv R(\theta)$，$\vv F(\theta)$ & 平面上转角 $\theta$ 的旋转矩阵；关于倾角 $\theta/2$ 的直线的反射矩阵 \\
		$\vv H$ & Householder 矩阵 $\vv I-2\vv u\vv u^\trans$ \\
		$\vv A\otimes\vv B$，$\vv A\circ\vv B$ & Kronecker 积；逐元素乘积（仅见于 C 组习题） \\
		\bottomrule
	\end{longtable}
\end{center}
\newpage
\textbf{空间与映射}

\begin{center}
	\begin{longtable}{l>{\raggedright\arraybackslash}p{0.6\linewidth}}
		\toprule
		记号 & 含义 \\
		\midrule
		\endhead
		$\spn(\vv\alpha_1,\ldots,\vv\alpha_s)$ & 向量组张成的子空间 \\
		$\dim V$，$\dim_{\mathbb{F}}V$ & 维数（需要强调数域时写出 $\mathbb{F}$） \\
		$W_1+W_2$，$W_1\oplus W_2$ & 子空间的和、直和 \\
		$V\times W$ & 向量空间的直积 \\
		$V\cong U$ & 同构 \\
		$\mathscr T,\ \mathscr S,\ \mathscr D$ & 线性映射与线性变换用花体字母；$\mathscr I$ 是恒等变换，$\mathscr O$ 是零变换 \\
		$\ker\mathscr T$，$\imm\mathscr T$ & 核、像；对矩阵 $\vv A$，$\ker\vv A$、$\imm\vv A$ 指 $\vv x\mapsto\vv A\vv x$ 的核与像 \\
		$\mathscr T|_W$ & $\mathscr T$ 在不变子空间 $W$ 上的限制 \\
		$L(V,W)$ & $V$ 到 $W$ 的全体线性映射 \\
		$V_\lambda$ & 属于特征值 $\lambda$ 的特征子空间 \\
		$\langle\vv\alpha,\vv\beta\rangle$ & 内积；$\mathbb{R}^n$ 中默认为标准内积 $\vv\alpha^\trans\vv\beta$ \\
		$\|\vv\alpha\|$ & 范数（长度） \\
		$\vv\alpha\perp\vv\beta$ & 正交 \\
		$W^\perp$ & 正交补 \\
		$\vv\alpha_W$ & $\vv\alpha$ 在子空间 $W$ 上的正交投影 \\
		$\vv x_0+W$ & 线性流形（仿射子空间） \\
		$V^*$，$W^0$，$\mathscr T^*$ & 对偶空间、零化子、对偶映射（第十四章） \\
		$\hat{\vv\alpha}$ & $\vv\alpha$ 在 $V^{**}$ 中的像：$f\mapsto f(\vv\alpha)$（第十四章） \\
		$\vv x\ge\zero$ & 向量的每个分量都非负 \\
		\bottomrule
	\end{longtable}
\end{center}

\begin{note}
	同一个字母在不同章节中可能另有所指，例如 $\vv P$ 既表示过渡矩阵，也表示转移矩阵与投影矩阵；$\vv\Sigma$ 在第十一章是奇异值构成的矩阵，在第十三章是协方差矩阵. 这时以该章的说明为准.
\end{note}
.
%% 依赖：导言区定义的 \vv \zero \trans \spn \rk \imm \diag \tr \hB \hC.
%% =====================================================================

\chapter*{本书记号约定}
\addcontentsline{toc}{chapter}{本书记号约定}
\markboth{本书记号约定}{本书记号约定}

\paragraph*{关于书名.} 在我国高校的课程体系中，``高等代数''通常指数学类专业为期两个学期的基础课，除线性代数外还包括较系统的多项式理论和空间理论；面向非数学专业，如经济、管理、法学等专业开设的同类课程，则多称``线性代数''，以方程组、矩阵、向量空间、特征值与二次型为主. 本书不带标记的核心内容与后者的教学要求相当，按下文的 48 学时方案讲授即可. 书中另外收入了多项式、对偶空间与双线性型等传统上属于高等代数的内容，供 72、96 学时方案和有兴趣的读者选用，本书书名由此而来.

\section*{一、学时标记与教学方案}

本书内容明显多于一个学期所能讲授的分量. 为方便不同课时的课程选用，本书使用下列记号进行区分：

\begin{center}
	\begin{tabular}{cl}
		\toprule
		标记 & 含义 \\
		\midrule
		（无标记） & 核心内容，48 学时方案即讲授，72、96 学时方案也都讲授 \\
		\textdagger & 72 学时方案起讲授（96 学时方案也讲授） \\
		\textdaggerdbl & 仅在 96 学时方案中讲授 \\
		* & 选读内容，任何方案都不作要求，供有兴趣的读者自学 \\
		\bottomrule
	\end{tabular}
\end{center}

一节标题带有标记时，节内未另加标记的小节与它同属一个方案；小节自己的标记优先. 各章的建议学时如下表. 表中``—''表示该方案不讲这一章.

\begin{center}
	\begin{tabular}{llccc}
		\toprule
		章 & 内容 & 48 学时 & 72 学时 & 96 学时 \\
		\midrule
		第一章 & 线性方程组与消元法 & 3 & 3 & 3 \\
		第二章 & 向量空间 & 5 & 6 & 7 \\
		第三章 & 矩阵 & 6 & 7 & 8 \\
		第四章 & 线性空间与线性变换 & 5 & 7 & 9 \\
		第五章 & 行列式 & 4 & 5 & 6 \\
		第六章 & 线性方程组解的结构 & 4 & 5 & 7 \\
		第七章 & 多项式 & — & 4 & 6 \\
		第八章 & 特征值与特征向量 & 6 & 7 & 9 \\
		第九章 & $\lambda$-矩阵 & — & 4 & 6 \\
		第十章 & 内积空间 & 3 & 6 & 7 \\
		第十一章 & 矩阵分解 & 3 & 5 & 7 \\
		第十二章 & 酉空间 & — & — & 6 \\
		第十三章 & 二次型 & 5 & 7 & 8 \\
		第十四章 & *对偶空间与双线性型 & — & — & — \\
		\midrule
		& 复习与机动 & 4 & 6 & 7 \\
		& 合计 & 48 & 72 & 96 \\
		\bottomrule
	\end{tabular}
\end{center}

\paragraph*{使用说明.}
\begin{itemize}
	\item \textbf{48 学时方案}讲授各章不带标记的内容，依次覆盖：向量与矩阵运算、消元法；线性相关与秩；矩阵的运算与特殊矩阵；线性空间与线性映射；行列式；矩阵的秩；线性方程组的解；特征值与相似对角化；内积与欧氏空间；二次型. 
	\item 48 学时方案不讲第七章. 第八章用到的代数基本定理与 Vieta 定理，以不加证明的形式收在附录~\ref{app:complex}，所需的复数知识也在那里. 不带标记的内容若引用了带标记小节中的结论，48 学时方案只要求了解结论.
	\item \textbf{72 学时方案}讲授不带标记与带 \textdagger 的内容，大致相当于一学期、每周 4—5 学时. 第七章讲``一元多项式环''``带余除法''``根''``代数基本定理''与``最小多项式''；第八章不讲 Cayley–Hamilton 定理一节，因此对角化判据以定理~\ref{thm:diagcriterion} 为准；第九章 $\lambda$-矩阵整章讲授.
	\item \textbf{96 学时方案}讲完全部无星号的内容，包括第十二章酉空间，适合两个学期、每学期 48 学时的课程. 第十四章与各章星号小节可作为讨论班或期末报告的选题.
	\item 三个部分末尾的\textbf{综合项目}不占课内学时. 建议 48 学时方案选做第一部分的一个项目；72 学时方案在第一、二部分中各选一个；96 学时方案三个部分各选一个. 评分量规见``第一部分综合项目''.
	\item 章末习题的 A 组覆盖各方案的核心内容（涉及 \textdagger、\textdaggerdbl 小节的题目，未讲授时可以跳过）. A 组的答案、B 组的提示见附录.
\end{itemize}


\section*{二、栏目}

除了定义、定理与证明，各章还穿插着六类固定栏目，用不同颜色的方框标出. 它们各有分工，读者可以按需要取舍，但不宜一概跳过.

各章开头依次给出本章地位、建议学时、学习目标、前置知识与学习提示. 学习目标中注明``72 学时方案起''或``96 学时方案''的条目，对应正文中带 \textdagger 或 \textdaggerdbl 的内容；其余条目属于各方案共同的要求.

\textbf{思维动机}是正文最重要的补充. 它回答了另一个重要问题，即``为什么会想到这么做''：一个定义为什么这样写，一个证明的关键一步从何而来，一个结论在全书的位置在哪里. 初读时若只能挑一类栏目来读，编者推荐多留意这一类.

\textbf{专业连线}与\textbf{建模实验}即为数学建模的相关内容. 前者篇幅较短，指出一个概念在专业里的对应物，例如价目表是向量、投入产出表是矩阵；后者是完整的案例，依次走过``问题—假设—模型—求解—解读—局限''六步. 其中最后一步最容易被忽略，也最值得用心：每个模型都建立在简化之上，知道它在哪里失效，和会用它同样重要. 投入产出、物流网络、判例引证、法条检索等几个例子会跨章节反复出现，请留意这些问题是怎样随着工具的更新而被研究得更透彻的.

\textbf{历史回响}讲述概念的来历：谁在什么问题的推动下提出了它，又经过了怎样的曲折. 栏目中的史实均注明出处，见书末参考文献. \textbf{常见误区}收集初学者最容易犯的错误，多数以反例的形式给出，适合在做习题之前进行阅读.

\textbf{计算实践}是使用 Python 进行实践的栏目. 栏目中会布置数个任务供读者进行探究，并附检查点用于检验结果. 除此以外，本栏目还有介绍实际的数值计算和计算机求解中的算法的作用.

\section*{三、编号与引用}
\begin{itemize}
	\item 定义、定理、引理、推论、命题、例题按节编号，例如``定理 2.4.3''是第二章第四节的第 3 个定理.
	\item 章末习题分为 A、B、C三组，每组各自编号. 正文中``第三章章末习题 B16''指第三章章末习题 B 组第 16 题. 正文各节中零星的``练习''按出现顺序编号，与章末习题不同.
	\item 证明结束用 $\blacksquare$ 标出.
\end{itemize}

\section*{四、本书数学符号一览}

\textbf{数与集合}

\begin{center}
	\begin{longtable}{l>{\raggedright\arraybackslash}p{0.6\linewidth}}
		\toprule
		记号 & 含义 \\
		\midrule
		\endhead
		$\mathbb{N},\ \mathbb{Z},\ \mathbb{Q},\ \mathbb{R},\ \mathbb{C}$ & 自然数集、整数集、有理数集、实数集、复数集 \\
		$\mathbb{F}$ & 一般的数域 \\
		$A \subset B$ & $A$ 是 $B$ 的子集，等价于 $B \supset A$ \\
		$\mathrm{i}$ & 虚数单位，$\mathrm i^2=-1$ \\
		$\bar z$ & 复数 $z$ 的共轭 \\
		$|z|$ & 复数 $z$ 的模 \\
		$\mathrm e^{\mathrm i\theta}$ & $\cos\theta+\mathrm i\sin\theta$（见附录``复数''一节） \\
		$\omega$ & 三次单位根 $\mathrm e^{2\pi\mathrm i/3}$ \\
		$\mathbb{F}_2$ & （取决于上下文）只含 $0,1$、按模 $2$ 运算的数系 \\
		$\sum_{i=1}^n a_i$，$\prod_{i=1}^na_i$ & 求和、求积（见附录``预备知识''） \\
		$\binom nk$ & 组合数 \\
		$\deg f$ & 多项式 $f$ 的次数 \\
		$g\mid f$ & 多项式（或整数）$g$ 整除 $f$ \\
		$\gcd(f,g)$ & 首一的最大公因式 \\
		$f:A\to B$，$x\mapsto f(x)$ & 映射及其对应法则 \\
		$g\circ f$ & 映射的复合，先 $f$ 后 $g$ \\
		$\mathbb{F}[x]$，$\mathbb{F}[x]_n$ & $\mathbb{F}$ 上的一元多项式全体；次数小于 $n$ 的多项式（含零多项式）全体 \\
		$C[a,b]$ & 闭区间 $[a,b]$ 上的实值连续函数全体 \\
		\bottomrule
	\end{longtable}
\end{center}
\newpage
\textbf{向量与矩阵}

\begin{center}
	\begin{longtable}{l>{\raggedright\arraybackslash}p{0.6\linewidth}}
		\toprule
		记号 & 含义 \\
		\midrule
		\endhead
		$\vv\alpha,\ \vv\beta,\ \vv x,\ \vv A$ & 向量与矩阵一律用加粗斜体 \\
		$\zero$ & 零向量 \\
		$\vv O$ & 零矩阵 \\
		$\vv I$，$\vv I_n$ & 单位矩阵、$n$ 阶单位矩阵 \\
		$\vv e_i$ & 标准单位向量：第 $i$ 个分量为 $1$，其余为 $0$ \\
		$\vv 1$ & 分量全为 $1$ 的列向量 \\
		$\mathbb{F}^n$ & $\mathbb{F}$ 上的 $n$ 维列向量全体；本书``向量''默认指列向量 \\
		$M_{m\times n}(\mathbb{F})$，$M_n(\mathbb{F})$ & $m\times n$ 矩阵全体；$n$ 阶方阵全体 \\
		$a_{ij}$，$(a_{ij})_{m\times n}$ & 矩阵的 $(i,j)$ 元；以 $a_{ij}$ 为元素的矩阵 \\
		$\vv A^\trans$ & 转置 \\
		$\vv A^{-1}$ & 逆矩阵 \\
		$\vv A^*$ & 伴随矩阵（第五章）；第十四章中 $\mathscr T^*$ 另指对偶映射 \\
		$(\vv A\mid\vv b)$ & 增广矩阵 \\
		$\diag(\lambda_1,\ldots,\lambda_n)$ & 对角矩阵 \\
		$r_i\leftrightarrow r_j$，$kr_i$，$r_i+cr_j$ & 三类初等行变换（列变换用 $c$ 代替 $r$） \\
		$\vv E(i,j)$，$\vv E(i(k))$，$\vv E(i,j(c))$ & 三类初等矩阵 \\
		$\rk\vv A$ & 秩 \\
		$\det\vv A$，$|\vv A|$ & 行列式 \\
		$\tr\vv A$ & 迹，即主对角元之和 \\
		$M_{ij}$，$A_{ij}$ & 余子式、代数余子式 \\
		$\tau(i_1i_2\cdots i_n)$ & 排列的逆序数 \\
		$f_{\vv A}(\lambda)$，$m_{\vv A}(\lambda)$ & 特征多项式 $\det(\lambda\vv I-\vv A)$，最小多项式 \\
		$\vv J_k(\lambda)$ & $k$ 阶 Jordan 块 \\
		$\sigma_1\ge\sigma_2\ge\cdots$ & 奇异值 \\
		$\kappa(\vv A)$ & 条件数 $\sigma_1/\sigma_n$ \\
		$\|\vv A\|_F$ & Frobenius 范数 $\sqrt{\sum a_{ij}^2}$ \\
				$\|\vv A\|_2$ & 谱范数，即最大奇异值 $\sigma_1$ \\
		$\vv A^+$ & 伪逆 \\
		$\mathrm e^{t\vv A}$ & 矩阵指数 \\
		$\vv R(\theta)$，$\vv F(\theta)$ & 平面上转角 $\theta$ 的旋转矩阵；关于倾角 $\theta/2$ 的直线的反射矩阵 \\
		$\vv H$ & Householder 矩阵 $\vv I-2\vv u\vv u^\trans$ \\
		$\vv A\otimes\vv B$，$\vv A\circ\vv B$ & Kronecker 积；逐元素乘积（仅见于 C 组习题） \\
		\bottomrule
	\end{longtable}
\end{center}
\newpage
\textbf{空间与映射}

\begin{center}
	\begin{longtable}{l>{\raggedright\arraybackslash}p{0.6\linewidth}}
		\toprule
		记号 & 含义 \\
		\midrule
		\endhead
		$\spn(\vv\alpha_1,\ldots,\vv\alpha_s)$ & 向量组张成的子空间 \\
		$\dim V$，$\dim_{\mathbb{F}}V$ & 维数（需要强调数域时写出 $\mathbb{F}$） \\
		$W_1+W_2$，$W_1\oplus W_2$ & 子空间的和、直和 \\
		$V\times W$ & 向量空间的直积 \\
		$V\cong U$ & 同构 \\
		$\mathscr T,\ \mathscr S,\ \mathscr D$ & 线性映射与线性变换用花体字母；$\mathscr I$ 是恒等变换，$\mathscr O$ 是零变换 \\
		$\ker\mathscr T$，$\imm\mathscr T$ & 核、像；对矩阵 $\vv A$，$\ker\vv A$、$\imm\vv A$ 指 $\vv x\mapsto\vv A\vv x$ 的核与像 \\
		$\mathscr T|_W$ & $\mathscr T$ 在不变子空间 $W$ 上的限制 \\
		$L(V,W)$ & $V$ 到 $W$ 的全体线性映射 \\
		$V_\lambda$ & 属于特征值 $\lambda$ 的特征子空间 \\
		$\langle\vv\alpha,\vv\beta\rangle$ & 内积；$\mathbb{R}^n$ 中默认为标准内积 $\vv\alpha^\trans\vv\beta$ \\
		$\|\vv\alpha\|$ & 范数（长度） \\
		$\vv\alpha\perp\vv\beta$ & 正交 \\
		$W^\perp$ & 正交补 \\
		$\vv\alpha_W$ & $\vv\alpha$ 在子空间 $W$ 上的正交投影 \\
		$\vv x_0+W$ & 线性流形（仿射子空间） \\
		$V^*$，$W^0$，$\mathscr T^*$ & 对偶空间、零化子、对偶映射（第十四章） \\
		$\hat{\vv\alpha}$ & $\vv\alpha$ 在 $V^{**}$ 中的像：$f\mapsto f(\vv\alpha)$（第十四章） \\
		$\vv x\ge\zero$ & 向量的每个分量都非负 \\
		\bottomrule
	\end{longtable}
\end{center}

\begin{note}
	同一个字母在不同章节中可能另有所指，例如 $\vv P$ 既表示过渡矩阵，也表示转移矩阵与投影矩阵；$\vv\Sigma$ 在第十一章是奇异值构成的矩阵，在第十三章是协方差矩阵. 这时以该章的说明为准.
\end{note}
